How similar are the hidden states at different layers of a multi-layer perceptron (MLP)? Prior work has tabulated cross-layer representation similarity for CNNs, ResNets, and transformers, but pure MLPs remain under-studied, even though their lack of attention, convolutions, and skip connections makes them the cleanest architecture in which to isolate the interplay of width, depth, and nonlinearity. We train a grid of ReLU MLPs on \cifar and \mnist with depths $\{4, 8, 16\}$ and widths $\{128, 512, 2048\}$ across three seeds, and measure cross-layer similarity with \ckam, orthogonal-Procrustes, sample-wise cosine, PC-1 variance ratios, and per-layer linear probes. Three findings emerge. First, MLP hidden states are neither identical nor unrelated: mean off-diagonal CKA sits firmly inside $(0,1)$ (CIFAR: 0.50--0.78, MNIST: 0.77--0.93), with adjacent layers more similar than far-apart layers ($p\!\ll\!10^{-3}$). Second, at depth 16 a large block of layers (2--14) collapses to CKA\,$\gtrsim\!0.9$, reproducing the block-structure phenomenon previously documented only for CNNs and transformers in a pure-MLP setting. Third, and contrary to the CNN literature, \emph{wider MLPs produce less similar layers} at depths 4 and 8 (Spearman $\rho\!=\!-0.95$, $p\!<\!10^{-3}$), a direct signature of limited per-layer transformation capacity; this reverses at depth 16, where block structure dominates. Three metrics agree on the ranking of layer pairs ($\rho\!=\!0.98$ between Procrustes and CKA), and controls (shuffled labels, lazy regime, dominant-datapoint ablation) rule out trivial explanations.
